\documentclass[11pt,a4paper]{article}
\usepackage[margin=24mm]{geometry}
\usepackage[UTF8,fontset=fandol,scheme=plain]{ctex}
\usepackage{amsmath,amssymb,amsthm,booktabs,array}
\usepackage[colorlinks=true,linkcolor=blue,urlcolor=blue,citecolor=blue]{hyperref}
\usepackage{xurl}
\newcommand{\lean}[1]{\nolinkurl{#1}}
\newcommand{\code}[1]{\texttt{\detokenize{#1}}}
\newcommand{\F}{\mathbb F}
\newcommand{\fpt}{\operatorname{fpt}}
\newcommand{\origin}{\mathfrak m_0}
\renewcommand{\refname}{References}
\setlength{\parindent}{0pt}
\setlength{\parskip}{5pt}
\setlength{\emergencystretch}{2em}
\title{Disproof of conjecture 00000000525\\[5pt]
\large An all-level Frobenius bound in characteristic five}
\author{C0ldSmi1e}
\date{5 October 2026}
\begin{document}
\maketitle
\begin{abstract}
For $f=x^2+y^3$ over $\F_5$, we prove
$\fpt_0(f)\le 4/5<5/6$ and the same upper bound for the global threshold.
Since $5$ is prime and $5\equiv2\pmod3$, this disproves the exact-value
assertion, and hence the conjunction in the conjecture. We do not assert
an exact threshold or decide its separate limiting assertion.
\end{abstract}

\section{Original statement and interpretation}
\textbf{English.} Definition: The F-pure threshold $\fpt(f)$ (the
characteristic-$p$ logarithmic threshold). Conjecture: The limit of
$\fpt(x^2+y^3)$ as $p\to\infty$ equals its characteristic-zero log
threshold, and for $p\equiv2\pmod3$, $\fpt=5/6$ (an explicit interpolation
formula).

\textbf{中文。} 定义：F-纯阈值 $\fpt(f)$（特征 $p$ 的 log 阈值）。
猜想：$\fpt(x^2+y^3)$ 的极限（$p\to\infty$）等于其特征零 log 阈值，
且对 $p\equiv2\pmod3$ 有 $\fpt=5/6$（显式插值公式）。

Neither version excludes small primes or replaces the exact-value assertion
by an eventual assertion. The witness $p=5$ also exceeds both exponents
in the polynomial. Both the threshold at the origin and the ordinary
global interpretation are addressed below.

For a field $K$ of characteristic $p$, put $R=K[x,y]$ and
$\origin=(x,y)$. We use the standard all-level definition\,\cite{sv}
\[
 \fpt_0(f)=\sup\left\{\frac{n}{p^e}: e\ge1,\ n\in\mathbb N,
 f^n\notin(x^{p^e},y^{p^e})\right\}.
\]
The denominator and supremum are in $\mathbb R$. The ideal here is the
Frobenius power $\origin^{[p^e]}$, not the ordinary ideal power.
Lean uses the actual type \lean{MvPolynomial (Fin 2) K} and
\lean{Ideal.span}; no numerical surrogate for either object is used.

\newpage
\section{A bound at every Frobenius level}
The following polynomial identity holds over every field:
\begin{align*}
 (x^2+y^3)^4
 &=x^5(x^3+4xy^3)\\
 &\quad+y^5(6x^4y+4x^2y^4+y^7).
\end{align*}
Consequently $f^4\in(x^5,y^5)$. This is checked as a polynomial identity
by Lean's kernel-verified ring normalization, followed by ordinary ideal
membership operations.

Now suppose $\operatorname{char}K=5$. For every $r\in\mathbb N$, the
$5^r$th power map is a ring homomorphism. Taking that power of an ideal
membership expression gives
\[
 f^{4\cdot5^r}\in
 (x^{5^{r+1}},y^{5^{r+1}}).
 \tag{1}
\]
More generally, the checked lemma \lean{frobenius_mem} sends membership
in $(x^q,y^q)$ to membership of the $5^r$th power in
$(x^{q5^r},y^{q5^r})$. Thus (1) holds for every level, with no finite
cutoff or experimental extrapolation.

An ideal containing $f^a$ contains $f^n$ whenever $n\ge a$.
Therefore, for any positive level $e=r+1$, nonmembership of $f^n$
forces $n<4\cdot5^r$. Every defining ratio then satisfies
\[
 \frac{n}{5^{r+1}}\le\frac{4\cdot5^r}{5^{r+1}}=\frac45.
\]
The defining set is nonempty: $n=0$ contributes zero, since $1$ is
outside $(x^{5^e},y^{5^e})$. Evaluating at the origin proves this
properness. The bound above proves boundedness. Conditional completeness
of $\mathbb R$ therefore yields the genuine supremum bound
\[
 \boxed{\fpt_0(x^2+y^3)\le\frac45<\frac56.}
 \tag{2}
\]

\paragraph{Indexing in the Lean definition.}
The implementation permits $e=0$. At that level the ideal is $(x,y)$,
and $f\in(x,y)$, so the only admissible exponent is $n=0$. Its ratio
zero is already present at level one. The theorem
\lean{allLevels_eq_positive} proves equality of the two defining sets;
\lean{fpt_eq_positive_levels} specializes it to this polynomial.
Thus the added zeroth level changes neither the set nor its supremum.
No empty-set or unbounded-set default of the real supremum is used.

\newpage
\section{The global interpretation}
For an ideal $I\subset R$ and a positive integer $q$, write
\[
 I^{[q]}=(g^q:g\in I).
\]
For every maximal ideal $I$ containing $f$, define
\[
 c^I(f)=\sup\left\{\frac{n}{p^e}: e\ge1,\ n\in\mathbb N,
 f^n\notin I^{[p^e]}\right\}.
\]
Over an F-finite regular ring, the standard global characterization is
$\fpt(f)=\inf_{I\supset(f),\ I\text{ maximal}}c^I(f)$;
the local interpretation and minimum characterization are given in
\cite[\S3, Proposition 3.2(iii--iv)]{bms}.
These hypotheses hold for the prime-field polynomial rings $\F_p[x,y]$
used in the final statement. Ideals are not restricted to rational points.
The equivalence with the test-ideal definition is cited; this project
formalizes the displayed characterization directly.

Lean formalizes the displayed bracket ideal, each supremum, and this
infimum directly. For $f\in I$ and $p>0$, one has
$f^{p^e}\in I^{[p^e]}$. Hence every admissible exponent satisfies
$n<p^e$, so each defining set is bounded above by one. It contains zero
because $I^{[p^e]}\subseteq I\ne R$. All its ratios are nonnegative.
The corresponding facts and the positive-level indexing bridge are
checked in \lean{thresholdAt_admissible_le_one},
\lean{thresholdAt_set_nonempty}, and
\lean{thresholdAt_eq_positive_levels}.

Let $\origin$ now be the kernel of constant-coefficient evaluation
$R\to K$. This map is surjective, so its kernel is maximal, and $f$
belongs to it. For every $q$,
\[
 (x^q,y^q)\subseteq\origin^{[q]}.
\]
Thus the nonmembership ratios for $\origin^{[5^e]}$ form a subset of
those defining $\fpt_0(f)$. Their suprema satisfy
$c^{\origin}(f)\le\fpt_0(f)$. The set defining the global infimum is
nonempty (it includes the origin term) and bounded below by zero.
Consequently Lean proves the exact chain needed:
\[
 \fpt(f)\le c^{\origin}(f)\le\fpt_0(f)\le\frac45<\frac56.
 \tag{3}
\]
The middle comparison uses the established nonemptiness and boundedness,
not a formal default value of an extremum. The combined theorem is
\lean{globalFpt_le_origin_fpt}. Equality of local and global thresholds
is neither asserted nor required.

The algebraic inequalities are proved for the displayed expressions over
every characteristic-five field. The cited identification as a global
F-pure threshold is used here for $\F_p[x,y]$, within the F-finite
regular setting; no assertion about arbitrary non-F-finite fields is
deduced from that reference.

\newpage
\section{Faithfulness and verification}
\lean{PrimeCharacteristic} is the subtype of all prime natural numbers.
\lean{localPrimeThreshold} and \lean{globalPrimeThreshold} use
$\mathbb Z/p\mathbb Z$, with its field instance obtained from primality.
\lean{ResidueClaim} states the equality $T(p)=5/6$ for every prime
$p$ with $p\bmod3=2$. Applying it at $5$ contradicts (2) or (3).

\lean{ConjectureAt T L} is the conjunction of that residue assertion
and convergence of $T(p)$ to $L$ through all primes tending to infinity.
The formal limit uses the pullback of the natural-number at-top filter
along the prime subtype's value map. The final theorems are
\begin{center}
\lean{Conjecture525.local_conjecture_false}\\[3pt]
\lean{Conjecture525.global_conjecture_false}.
\end{center}
Each proves the negation for \emph{every} real $L$. In particular, it
applies to the actual characteristic-zero log canonical threshold,
whatever its value. $L$ is not a redefinition of that invariant, and the
proof makes no assumption that the separate limit assertion fails.

\paragraph{Reproducibility.}
The project pins Lean 4.19.0 and Mathlib v4.19.0, including all eight
transitive dependency revisions in \code{lake-manifest.json}.
Run \code{lake build} in \code{lean/}; then replay
\code{Conjecture525.lean}, \code{Audit.lean}, and \code{Check.lean}
with \code{lake env lean -DwarningAsError=true}.
The accompanying README gives complete commands and artifact locations.

The author and a separate coordinator build both completed successfully.
The coordinator used a fresh project output directory and recorded 45
commands, including dependency identities before and after execution.
All three modules passed strict replay. Independent source and compiled
environment inventories cover 48 named authored declarations (15
definitions/abbreviations and 33 lemmas/theorems) and all 15 additional
generated declarations: 63 in total. Every compiled declaration is safe,
with no partial or missing dependency, and no axiom declaration. The only
axiom dependencies are the standard foundations \code{propext},
\code{Classical.choice}, and \code{Quot.sound}. No admission,
\code{native_decide}, extra axiom, or compiler-artifact exemption is used.

\code{Audit.lean} and \code{Check.lean} inspect existing declarations;
they do not generate a mathematical theorem. All mathematical
calculations are in Lean. There is no auxiliary numerical computation
whose output is assumed as a proof. Evidence preserves actual successful
and failed development commands and clearly separates local verification
from any later maintainer decision.

\begin{thebibliography}{9}
\bibitem{sv} K. E. Smith and A. Vraciu, \emph{Values of the F-pure
threshold for homogeneous polynomials}, J. London Math. Soc. 108 (2023),
1004--1035, Definition 2.1.
\href{https://doi.org/10.1112/jlms.12774}{doi:10.1112/jlms.12774}.
Only the general definition is used, not homogeneous-polynomial formulas.
\bibitem{bms} M. Blickle, M. Musta\c{t}\u{a}, and K. E. Smith,
\emph{F-thresholds of hypersurfaces}, \S3 and Proposition 3.2,
\href{https://arxiv.org/abs/0705.1210}{arXiv:0705.1210v3} (2008).
\end{thebibliography}
\end{document}
